% Part: computability
% Chapter: computability-theory
% Section: ce-closed-cup-cap

\documentclass[../../../include/open-logic-section]{subfiles}

\begin{document}

\olfileid{cmp}{thy}{clo}

\olsection[Gabungan lan Irisan Himpunan C.E.]{Himpunan kang Bisa Dienumerasi
  sacara Komputabel Katutup tumrap Gabungan lan Irisan}

Teorema ing ngisor iki menehi sawetara sipat
katutupan kanggo himpunan kang bisa dienumerasi
sacara komputabel.

\begin{thm}
Upamane $A$ lan $B$ bisa dienumerasi sacara
komputabel. Mula $A \cap B$ lan $A \cup B$
uga bisa dienumerasi sacara komputabel.
\end{thm}

\begin{proof}
\olref[eqc]{thm:ce-equiv} ngidini kita
nggunakake maneka ciri himpunan kang bisa
dienumerasi sacara komputabel. Kanggo
nggambarake iki, kita bakal menehi sawetara
bukti kang beda.

Kanggo bukti kapisan, upamane $A$ dienumerasi
dening fungsi komputabel~$f$, lan $B$
dienumerasi dening fungsi komputabel~$g$.
Tetepna
\begin{align*}
h(x) & = \umin{y}{(f(y) = x \lor g(y) = x)} \text{ lan}\\
j(x) & = \umin{y}{(f((y)_0) = x \land g((y)_1) = x)}.
\end{align*}
Mula $A \cup B$ minangka domain $h$,
lan $A \cap B$ minangka domain~$j$.

\begin{explain}
Mangkene teges komputasionale: yen diwenehi
prosedur kang ngenumerasi $A$ lan $B$,
kita bisa semi-mutusake apa unsur $x$
kalebu $A \cup B$ kanthi nggoleki $x$
ing salah siji enumerasi; lan kita bisa
semi-mutusake apa unsur $x$ kalebu
$A \cap B$ kanthi nggoleki $x$ ing
loro enumerasi bebarengan.
\end{explain}

Kanggo bukti kapindho, upamane maneh $A$
dienumerasi dening~$f$ lan $B$ dening~$g$.
Tetepna
\[
k(x) = \begin{cases}
f(x/2) & \text{yen $x$ genap} \\
g((x-1)/2) & \text{yen $x$ ganjil.}
\end{cases}
\]
Mula $k$ ngenumerasi $A \cup B$;
gagasane, $k$ mung nggenteni enumerasi
saka $f$ lan~$g$ kanthi selang-seling.
Ngenumerasi $A \cap B$ luwih angel.
Yen $A \cap B$ kosong, himpunan iki
cetha bisa dienumerasi sacara komputabel.
Yen ora, tetepna $c$ minangka unsur
sembarang saka $A \cap B$, banjur
tetepna $l$ minangka
\[
l(x) = \begin{cases}
f((x)_0) & \text{yen $f((x)_0) = g((x)_1)$} \\
c & \text{yen ora.}
\end{cases}
\]
Miturut komputasi, $l$ mriksa pasangan
unsur ing enumerasi $f$ lan $g$,
banjur ngasilake saben nilai kang
padha sing ditemokake; yen ora padha, fungsi
iki mung ngasilake $c$ maneh.

Kanggo bukti pungkasan, upamane $A$
minangka \emph{domain} fungsi parsial
$m(x)$ lan $B$ minangka domain fungsi
parsial $n(x)$. Mula $A \cap B$
minangka domain fungsi parsial $m(x) +
n(x)$.

\begin{explain}
Miturut komputasi, yen $A$ minangka
himpunan nilai kang ndadekake $m$ mandheg
lan $B$ minangka himpunan nilai kang
ndadekake $n$ mandheg, mula $A \cap B$
minangka himpunan nilai kang ndadekake
loro prosedur iku padha mandheg.
\end{explain}

Nuduhake $A \cup B$ minangka himpunan
nilai kang ndadekake komputasi mandheg
luwih angel, amarga $m$ lan $n$ kudu
disimulasi bebarengan. Tetepna $d$
minangka indeks kanggo $m$ lan $e$
minangka indeks kanggo $n$; kanthi
tembung liya, $m = \cfind{d}$ lan
$n = \cfind{e}$. Mula $A \cup B$
minangka domain fungsi
\[
p(x) = \umin{y}{(T(d,x,y) \lor T(e,x,y))}.
\]
\begin{explain}
Miturut komputasi, ing input $x$,
$p$ nggoleki komputasi kang mandheg
kanggo $m$ utawa kanggo $n$,
banjur mandheg yen nemokake salah siji.
\end{explain}
\end{proof}

\end{document}
